% Equations used in the 3Di-conditioned flow-matching decoder (Steps E--T).
%
% Per-variant training and sampling formulas. Compiles standalone:
%   pdflatex equations_3di_flow_decoder.tex
%
% Source-of-truth code references:
%   src/lobster/model/latent_generator/tokenizer/_tokenizer_3di_input_flow.py
%   src/lobster/model/latent_generator/structure_decoder/_vit_decoder_conditional.py
%   .venv/lib/python3.12/site-packages/bionemo/moco/interpolants/.../continuous_flow_matching.py

\documentclass[11pt]{article}

\usepackage[margin=1in]{geometry}
\usepackage{amsmath, amssymb, amsthm}
\usepackage{booktabs}
\usepackage{array}
\usepackage{xcolor}
\usepackage{hyperref}
\hypersetup{colorlinks=true, urlcolor=blue, linkcolor=black}

\title{Equations for the 3Di-Conditioned Flow-Matching Decoder}
\author{Lobster project (\texttt{lobster.model.latent\_generator.tokenizer.Tokenizer3diInputFlow})}
\date{Snapshot 2026-06-08 (Steps E--T)}

\begin{document}
\maketitle

\section*{Notation}

\begin{tabular}{ll}
\toprule
$x_1 \in \mathbb{R}^{B\times L \times A \times 3}$ & ground-truth backbone coords (\AA), with $A=3$ atoms (N, $C_\alpha$, C).\\
$x_0 \in \mathbb{R}^{B\times L \times A \times 3}$ & prior sample.\\
$x_t$ & interpolated state at time $t \in [0,1]$ ($t=0$ pure noise, $t=1$ data).\\
$m \in \{0,1\}^{B\times L}$ & residue mask (1=valid, 0=pad).\\
$s \in \{0,\dots,19,20\}^{B\times L}$ & 3Di tokens (20 classes + 1 null/pad/CFG-drop class).\\
$\sigma$ & coord scale ($\sigma=0.1$, \AA{}$\to$nm; Step J).\\
$f_\theta(\cdot)$ & conditional decoder (\texttt{ViTDecoderConditional}); outputs match the moco interpolant's prediction type.\\
$\hat{x}_1$ & data-space prediction; $\hat{v}$ & velocity-space prediction.\\
$p_{\text{drop}}$ & CFG dropout prob (default $0.15$).\\
$w$ & CFG guidance scale (default $1.0$ at training, $\geq 1$ at inference).\\
$\lambda_{\text{aux}}$ & aux pairwise weight (default $0.01$).\\
$\epsilon$ & numerical floor ($10^{-5}$ in moco); subscript $m$ on $\|\cdot\|_m$ = mask-weighted mean.\\
\bottomrule
\end{tabular}

\section{Common machinery (all variants)}

\subsection{Centering (geometric CoM over masked atoms)}
\begin{equation}
\operatorname{Center}(x)_{b,l,a} \;=\; m_{b,l}\!\left(x_{b,l,a} \;-\; \frac{\sum_{l',a'} m_{b,l'} \, x_{b,l',a'}}{A \cdot \sum_{l'} m_{b,l'}}\right).
\end{equation}

\subsection{Coordinate scaling (Step J)}
Applied once after centering; \emph{undone} at the end of \texttt{sample()}:
\begin{equation}
\tilde x_1 \;=\; \sigma \cdot \operatorname{Center}(x_1), \qquad \sigma = 0.1.
\end{equation}

\subsection{Prior and zero-CoM noise (Step E delta \#2)}
\begin{equation}
\bar x_0 \sim \mathcal{N}(0, I_{BLA\times 3}), \qquad x_0 \;=\; \operatorname{Center}(\bar x_0).
\end{equation}

\subsection{Time distribution}
\begin{equation}
t \;\sim\; \mathcal{U}(0,1), \quad \text{i.i.d.\ per sample.}
\end{equation}

\subsection{Linear interpolant (moco \texttt{ContinuousFlowMatcher})}
\begin{equation}
x_t \;=\; t\, \tilde x_1 \;+\; (1-t)\, x_0, \qquad t=1 \Rightarrow x_t=\tilde x_1, \quad t=0 \Rightarrow x_t=x_0.
\end{equation}

\subsection{CFG conditioning dropout (Step E delta, train-only)}
With probability $p_{\text{drop}}$ (default $0.15$), replace the 3Di token sequence with the null class $20$:
\begin{equation}
\mathbf{1}^{\text{full}}_b \;\sim\; \mathrm{Bernoulli}(p_{\text{drop}}), \qquad s'_b \;=\; \begin{cases} \mathbf{1}_{L} \cdot 20 & \text{if } \mathbf{1}^{\text{full}}_b = 1 \\ s_b & \text{otherwise.} \end{cases}
\end{equation}

\subsection{Per-residue 3Di masking (Step V, train-only, optional)}
When \texttt{mask\_3di\_per\_residue=True}, samples NOT hit by the full CFG drop ($\mathbf{1}^{\text{full}}_b = 0$) get an additional per-token corruption: sample a per-sample mask rate $p_b \sim \mathcal{U}(0, 1)$, then replace each residue's 3Di token by the null class independently with probability $p_b$:
\begin{align}
p_b &\;\sim\; \mathcal{U}(0, 1), \\
M_{b,l} &\;\sim\; \mathrm{Bernoulli}(p_b), \\
s''_{b,l} &\;=\; \begin{cases} 20 & \text{if } M_{b,l} = 1 \text{ and } \mathbf{1}^{\text{full}}_b = 0 \\ s'_{b,l} & \text{otherwise.} \end{cases}
\end{align}
The marginal distribution of the per-sample mask fraction $\frac{1}{L}\sum_l M_{b,l}$ is approximately $\mathcal{U}(0, 1)$ for large $L$, so the model is calibrated at every coverage level $\in [0, 1]$ -- not concentrated at a single rate (which a fixed-rate per-token Bernoulli would produce, with Binomial concentration around the rate). Same intuition as $t \sim \mathcal{U}(0, 1)$ in diffusion training. Pattern borrowed from moco's discrete flow matching ([\texttt{discrete\_flow\_matching.py:73}]) and leflur's epitope masking convention.

\subsection{Random SO(3) augmentation (Step H, train-only, optional)}
When \texttt{random\_so3\_aug=True}, sample $R_b \sim \operatorname{Uniform}(\mathrm{SO}(3))$ via the QR + sign trick (Mezzadri 2007) and rotate the centered, scaled GT:
\begin{equation}
\tilde x_1 \;\leftarrow\; \tilde x_1 \cdot R_b^\top, \quad R_b R_b^\top = I, \quad \det R_b = +1.
\end{equation}
Applied \emph{after} centering and \emph{before} Kabsch / interpolation, so the rotated $\tilde x_1$ is what the loss target consumes.

\subsection{Kabsch coupling (Step E delta, optional)}
When \texttt{augmentation\_type=kabsch}, align $x_0$ to $\tilde x_1$ per sample by the optimal rotation:
\begin{equation}
R^\star_b \;=\; \arg\min_{R \in \mathrm{SO}(3)} \, \|R\, x_{0,b} \;-\; \tilde x_{1,b}\|_m^2, \qquad x_0 \leftarrow R^\star \cdot x_0.
\end{equation}
\emph{Note:} when both \texttt{random\_so3\_aug} and Kabsch are on, the Kabsch step aligns to the rotated $\tilde x_1$, which absorbs the SO(3) augmentation (Step S analysis).

\section{Decoder}

The conditional U-ViT decoder consumes \emph{token embedding + noisy-coord projection} per residue:
\begin{equation}
e_{b,l} \;=\; E_{\text{tok}}(s'_{b,l}) \;+\; W_{\text{coord}} \, \mathrm{vec}(x_{t,b,l}) \;+\; \mathbb{1}\!\left[\text{self-cond}\right] \cdot W_{\text{sc}} \, \mathrm{vec}\!\left(x^{\text{sc}}_{b,l}\right),
\end{equation}
with optional learnable register tokens (Proteina-style attention sinks) prepended. Time is lifted by a sinusoidal \texttt{NoiseConditioningBlock} into the FiLM-gated U-ViT blocks. The output type is governed by the moco interpolant's \texttt{prediction\_type}:
\begin{equation}
\hat{y} \;=\; f_\theta(s', m, \text{idx}, x_t, t \,;\, x^{\text{sc}}) \;\in\;
\begin{cases}
\text{data space, } \hat{x}_1 & \text{if \texttt{prediction\_type=DATA},}\\
\text{velocity space, } \hat{v} & \text{if \texttt{prediction\_type=VELOCITY}.}
\end{cases}
\end{equation}

For self-conditioning (\texttt{use\_self\_conditioning=True}), the second projection $W_{\text{sc}}$ is \emph{zero-initialised}.

\subsection{Cross-space conversion (moco helpers)}

Both helpers dispatch on the interpolant's \texttt{prediction\_type}; for the ``other'' parametrisation they convert via the linear interpolant relation $\hat{x}_1 = x_t + (1-t)\,\hat{v}$:
\begin{align}
\operatorname{ToData}(\hat{y}; x_t, t) &\;=\; \begin{cases} \hat{x}_1 & \text{(DATA)} \\ x_t + (1-t)\,\hat{v} & \text{(VELOCITY)} \end{cases} \\[0.4em]
\operatorname{ToVel}(\hat{y}; x_t, t) &\;=\; \begin{cases} (\hat{x}_1 - x_t)/(1 - t + \epsilon) & \text{(DATA)} \\ \hat{v} & \text{(VELOCITY)} \end{cases}
\end{align}

\section{Training losses}

\subsection{Flow-matching loss (moco \texttt{loss(\dots, target\_type)})}

The general form is
\begin{equation}
\mathcal{L}_{\text{FM}} \;=\; \mathbb{E}_{x_1, x_0, t}\Big[\, w(t)\cdot \|\,\Pi(\hat{y}; x_t, t) \;-\; y_{\text{target}}\,\|_m^2 \,\Big],
\end{equation}
where $\Pi$ converts the model output to the target space and $w(t)$ is the moco-internal time weighting:
\begin{center}
\begin{tabular}{lll}
\toprule
\texttt{target\_type} & $y_{\text{target}}$ & $w(t)$ \\
\midrule
DATA & $\tilde x_1$ & $1 / ((1-t)^2 + \epsilon)$ \\
VELOCITY & $\tilde x_1 - x_0$ & $1$ \\
\bottomrule
\end{tabular}
\end{center}

In the lobster code, \texttt{target\_type} is read from the interpolant's \texttt{prediction\_type}, so the canonical pairing is target=prediction:

\paragraph{DATA prediction $\Rightarrow$ DATA target (used by \texttt{flow}, \texttt{flow\_base}, \texttt{flow\_selfcond}, \texttt{flow\_nokabsch})}
\begin{equation}
\boxed{\;
\mathcal{L}_{\text{DATA}} \;=\; \mathbb{E}_{x_1, x_0, t}\Big[\, \frac{1}{(1-t)^2 + \epsilon} \, \|\,\hat{x}_1 \;-\; \tilde x_1 \,\|_m^2 \,\Big].
\;}
\end{equation}

\paragraph{VELOCITY prediction $\Rightarrow$ VELOCITY target (used by \texttt{flow\_nokabsch\_velocity})}
\begin{equation}
\boxed{\;
\mathcal{L}_{\text{VEL}} \;=\; \mathbb{E}_{x_1, x_0, t}\Big[\, \|\,\hat{v} \;-\; (\tilde x_1 - x_0)\,\|_m^2 \,\Big].
\;}
\end{equation}
The $1/(1-t)^2$ reweight is dropped automatically (no $w(t)$ term for \texttt{target\_type=VELOCITY}).

\subsection{Auxiliary pairwise L2 (always computed in data space)}
\begin{equation}
\mathcal{L}_{\text{aux}} \;=\; \mathbb{E}\!\left[\,\frac{1}{\sum_{i,j} m_i m_j}\!\sum_{i,j} m_i m_j \big(\,d_{ij}(\hat{x}_{1}^\star)\;-\;d_{ij}(\tilde x_1)\,\big)^2 \,\right],
\end{equation}
with $\hat{x}_{1}^\star = \operatorname{ToData}(\hat{y}; x_t, t)$ and $d_{ij}(x) = \|x_{i, C_\alpha} - x_{j, C_\alpha}\|_2$.

\subsection{Total objective}
\begin{equation}
\mathcal{L}_{\text{total}} \;=\; \mathcal{L}_{\text{FM}} \;+\; \lambda_{\text{aux}} \, \mathcal{L}_{\text{aux}}, \qquad \lambda_{\text{aux}} = 0.01.
\end{equation}

\subsection{Self-conditioning warm forward (Steps F, U; for self-cond runs)}
With probability $p_{\text{sc}} = 0.5$ at training:
\begin{align}
\hat{y}_{\text{prev}} &\;=\; f_\theta(s', m, \text{idx}, x_t, t \,;\, x^{\text{sc}} = \mathbf{0}), \\
x^{\text{sc}} &\;=\; \mathrm{detach}\!\left(\, \operatorname{ToData}(\hat{y}_{\text{prev}};\, x_t,\, t) \,\right), \\
\hat{y} &\;=\; f_\theta(s', m, \text{idx}, x_t, t \,;\, x^{\text{sc}}),
\end{align}
otherwise $x^{\text{sc}} = \mathbf{0}$ (cold start). The aux projection $W_{\text{sc}}$ is zero-initialised so a fresh model is byte-identical to the non-self-cond variant on the first step.

\paragraph{Step U refinement.} The self-cond carry is the DATA-space estimate $\operatorname{ToData}(\hat{y}; x_t, t)$, not the raw $\hat{y}$. For DATA-pred models the conversion is a no-op so \texttt{flow\_selfcond} is byte-identical; for VELOCITY-pred models the conversion is $x_t + (1-t)\hat{v}$, i.e.\ ``what would $x_1$ look like if I straight-line integrated from $x_t$ at time $t$.'' Keeps the self-cond semantic (``feed the previous denoised estimate of $x_1$'') consistent across parametrisations.

\section{Sampling (inference)}

\subsection{Inference schedule}
Default linear, $N$ steps over $[0, 1)$:
\begin{equation}
t_k \;=\; \frac{k}{N}, \quad k = 0, 1, \ldots, N-1, \qquad \Delta t \;=\; \frac{1}{N}.
\end{equation}

\subsection{Initialisation}
\begin{equation}
x_{t_0} \;=\; \tau \cdot \operatorname{Center}(x_0), \qquad x_0 \sim \mathcal{N}(0, I), \quad \tau = \texttt{init\_temperature} \;\;(\text{default }1).
\end{equation}

\subsection{Conditional / unconditional / autoguidance branches}
At each step $k$ (and any $x^{\text{sc}}$ from step $k-1$):
\begin{align}
\hat{y}_{\text{cond}} &\;=\; f_\theta(s,\, m,\, \text{idx},\, x_{t_k},\, t_k \,;\, x^{\text{sc}}), \\
\hat{y}_{\text{null}} &\;=\; f_\theta(\mathbf{20},\, m,\, \text{idx},\, x_{t_k},\, t_k \,;\, x^{\text{sc}}) \quad \text{(used iff $w \neq 1$)}.
\end{align}

\subsection{CFG combine (default Proteina form)}
\begin{equation}
\hat{y}_{\text{used}} \;=\; \hat{y}_{\text{null}} \;+\; w \cdot \big(\,\hat{y}_{\text{cond}} \;-\; \hat{y}_{\text{null}}\,\big), \qquad w \geq 1.
\end{equation}

When an autoguidance model $f_{\text{ag}}$ is supplied with ratio $\rho$:
\begin{equation}
\hat{y}_{\text{used}} \;=\; w\,\hat{y}_{\text{cond}} \;+\; (1-w) \cdot \big(\,\rho\,\hat{y}_{\text{ag}} \;+\; (1-\rho)\,\hat{y}_{\text{null}}\,\big).
\end{equation}

\subsection{Vector field at the integrator}
moco's \texttt{step()} converts the model output to a velocity:
\begin{equation}
v(x_{t_k}, t_k) \;=\; \operatorname{ToVel}(\hat{y}_{\text{used}}; x_{t_k}, t_k) \;=\; \begin{cases} (\,\hat{x}_{1}^{\text{used}} - x_{t_k}\,) / (1 - t_k + \epsilon) & \text{(DATA model)} \\ \hat{v}^{\text{used}} & \text{(VELOCITY model)} \end{cases}
\end{equation}

\subsection{ODE Euler step}
\begin{equation}
\boxed{\; x_{t_{k+1}} \;=\; x_{t_k} \;+\; v(x_{t_k}, t_k)\,\Delta t,\qquad x_{t_{k+1}} \leftarrow \operatorname{Center}(x_{t_{k+1}}). \;}
\end{equation}
The re-centering step (always on in \texttt{sample()}) projects the field onto the zero-CoM subspace each step.

\subsection{SDE step (score-stochastic, optional)}
Vector-field $\to$ score conversion:
\begin{equation}
s(x_t, t) \;=\; \frac{t \, v(x_t, t) - x_t}{1 - t + \epsilon}.
\end{equation}
With moco's $g_t$ (mode \texttt{tan}/\texttt{us}/\texttt{1/t}) and temperatures $\tau_n, \tau_s$:
\begin{equation}
x_{t_{k+1}} \;=\; x_{t_k} \;+\; \big(\,v(x_{t_k}, t_k) \;+\; g_{t_k}\,\tau_s\,s(x_{t_k}, t_k)\,\big)\,\Delta t \;+\; \sqrt{2\,g_{t_k}\,\Delta t}\,\tau_n\,\eta_k,\quad \eta_k \sim \mathcal{N}(0, I),
\end{equation}
followed by re-centering. For $t > t_{\text{lim,ode}}$ the SDE step degrades to the ODE step (numerical safety near $t \to 1$).

\paragraph{Note (Step P).} SDE sampling requires \texttt{augmentation\_type=null} during training: moco forbids the vector-field $\to$ score conversion under OT/Kabsch coupling. So only the Step S runs (\texttt{flow\_nokabsch{,\_velocity}}) can run SDE; the older \texttt{flow*}/\texttt{flow\_base*} variants are ODE-only.

\subsection{Self-conditioning carry (for self-cond runs)}
Initialised $x^{\text{sc}}_0 = \mathbf{0}$. After each step (Step U: convert to DATA space \emph{using the pre-step $x_{t_k}$} so the carry's reference point matches the prediction's):
\begin{equation}
x^{\text{sc}}_{k+1} \;=\; \operatorname{detach}\!\left(\, \operatorname{ToData}\!\left(\,\hat{y}_{\text{used}}^{(k)};\; x_{t_k},\; t_k\,\right) \,\right).
\end{equation}
At $w=1$ this collapses to the conditional prediction (in data space); for $w>1$ the carry is the converted CFG-blended output (a slight train/sample mismatch, see Step S analysis). The conversion is a no-op for DATA-pred models (\texttt{flow\_selfcond}) and is $x_{t_k} + (1-t_k)\hat{v}$ for VELOCITY-pred models (\texttt{flow\_nokabsch\_velocity\_selfcond}).

\subsection{Output rescaling}
The integrator runs entirely in nm; the final output is rescaled back to \AA:
\begin{equation}
x_{1,\text{out}} \;=\; x_{t_N} \,/\, \sigma.
\end{equation}

\section{Per-variant summary}

The five in-flight runs differ only in the table below. Architecture, optimiser, time distribution, prior, schedule, $\sigma$, $\lambda_{\text{aux}}$, $p_{\text{drop}}$, $w_{\text{train}} = 1$, and aux loss are identical across runs (except where noted).

\renewcommand{\arraystretch}{1.15}
\begin{center}
\begin{tabular}{l c c c c l c c}
\toprule
\textbf{run} & \textbf{scale} & \textbf{prediction} & \textbf{Kabsch} & \textbf{SO(3) aug} & \textbf{FM loss} & \textbf{self-cond} & \textbf{3Di mask} \\
\midrule
\texttt{flow}                                  & S (20.6\,M)  & DATA     & on  & off & $\mathcal{L}_{\text{DATA}}$ & off & off \\
\texttt{flow\_base}                            & BASE (110\,M)& DATA     & on  & off & $\mathcal{L}_{\text{DATA}}$ & off & off \\
\texttt{flow\_selfcond}                        & S (20.6\,M)  & DATA     & on  & off & $\mathcal{L}_{\text{DATA}}$ & on  & off \\
\texttt{flow\_nokabsch}                        & S (20.6\,M)  & DATA     & off & on  & $\mathcal{L}_{\text{DATA}}$ & off & off \\
\texttt{flow\_nokabsch\_velocity}              & S (20.6\,M)  & VELOCITY & off & on  & $\mathcal{L}_{\text{VEL}}$  & off & off \\
\texttt{flow\_nokabsch\_velocity\_selfcond}    & S (20.6\,M)  & VELOCITY & off & on  & $\mathcal{L}_{\text{VEL}}$  & on  & off \\
\texttt{flow\_nokabsch\_velocity\_base}        & BASE (110\,M)& VELOCITY & off & on  & $\mathcal{L}_{\text{VEL}}$  & off & off \\
\texttt{flow\_nokabsch\_velocity\_mask3di}     & S (20.6\,M)  & VELOCITY & off & on  & $\mathcal{L}_{\text{VEL}}$  & off & on  \\
\bottomrule
\end{tabular}
\end{center}

\paragraph{What changes per row:}
\begin{itemize}\itemsep -1pt
\item \emph{scale}: $d_\text{model}$ / depth / heads / inner attn dim. S = $512/6/8/256$, BASE = $768/12/12/768$ (with \texttt{time\_cond\_dim}=$256$ at BASE vs $128$ at S).
\item \emph{prediction}: governs the moco interpolant's \texttt{prediction\_type} and (via the lobster code's read of that flag) the loss \texttt{target\_type}. Velocity rows use $\mathcal{L}_{\text{VEL}}$ \emph{without} the $1/(1-t)^2$ reweight.
\item \emph{Kabsch}: per-sample $x_0 \to R^\star x_0$ alignment in the training pipeline. With Kabsch on, SDE is forbidden by moco.
\item \emph{SO(3) aug}: random rotation of $\tilde x_1$ before interpolation. With Kabsch on, the alignment absorbs the rotation.
\item \emph{self-cond}: enables $W_{\text{sc}}$ on the decoder and the $p_{\text{sc}}=0.5$ warm-forward branch.
\end{itemize}

\paragraph{Training-pipeline pseudocode (one minibatch):}
\begin{enumerate}\itemsep -1pt
\item $\tilde x_1 \leftarrow \sigma\cdot\operatorname{Center}(x_1)$.
\item If \texttt{random\_so3\_aug}: $\tilde x_1 \leftarrow R\,\tilde x_1$, $R \sim \mathrm{SO}(3)$ uniform.
\item $x_0 \leftarrow \operatorname{Center}(\bar x_0)$, $\bar x_0 \sim \mathcal{N}(0, I)$.
\item If \texttt{augmentation\_type=kabsch}: $x_0 \leftarrow R^\star \, x_0$.
\item $t \sim \mathcal{U}(0,1)$; $x_t = t\,\tilde x_1 + (1-t)\,x_0$.
\item CFG dropout: $s' \leftarrow $ null token w.p.\ $p_{\text{drop}}$.
\item (Self-cond only) warm forward, detach.
\item $\hat{y} = f_\theta(s', m, \text{idx}, x_t, t \,;\, x^{\text{sc}})$.
\item Compute $\mathcal{L}_\text{FM}$ via the moco \texttt{loss(\dots, target\_type)} matching the interpolant's \texttt{prediction\_type}, and $\mathcal{L}_\text{aux}$ on $\operatorname{ToData}(\hat{y}; x_t, t)$.
\item $\mathcal{L}_\text{total} = \mathcal{L}_\text{FM} + \lambda_\text{aux}\,\mathcal{L}_\text{aux}$; backprop.
\end{enumerate}

\paragraph{Sampling-pipeline pseudocode (one rollout, $N$ steps):}
\begin{enumerate}\itemsep -1pt
\item $x_{t_0} \leftarrow \tau \cdot \operatorname{Center}(x_0)$, $x_0 \sim \mathcal{N}(0, I)$.
\item For $k = 0, 1, \dots, N-1$:
  \begin{enumerate}
  \item $\hat{y}_\text{cond} \leftarrow f_\theta(s,\,m,\,\text{idx},\,x_{t_k},\,t_k\,;\,x^\text{sc})$.
  \item If $w \neq 1$: $\hat{y}_\text{null} \leftarrow f_\theta(\mathbf{20},\,m,\,\text{idx},\,x_{t_k},\,t_k\,;\,x^\text{sc})$, then CFG-combine.
  \item $v_k \leftarrow \operatorname{ToVel}(\hat{y}_\text{used};\,x_{t_k},\,t_k)$ (no-op for VELOCITY model).
  \item ODE: $x_{t_{k+1}} = x_{t_k} + v_k\,\Delta t$. \quad SDE: as in Eq.~(20).
  \item Re-center: $x_{t_{k+1}} \leftarrow \operatorname{Center}(x_{t_{k+1}})$.
  \item (Self-cond only) $x^{\text{sc}} \leftarrow \operatorname{detach}(\hat{y}_\text{used})$.
  \end{enumerate}
\item Return $x_{1,\text{out}} = x_{t_N} / \sigma$ (back to \AA).
\end{enumerate}

\section{Why these choices (one-line rationale per delta)}

\begin{itemize}\itemsep -1pt
\item \emph{Coord scale $\sigma=0.1$ (Step J).} Centered backbone std $\approx 6$--$15$\,\AA{} $\gg 1$\,nm; without scaling the prior is negligible vs.\ data and the flow degenerates to a 1-shot regressor.
\item \emph{Zero-CoM noise + center every step (Step E delta \#2).} CoM-free flow; without it the model spends capacity learning translation invariance.
\item \emph{Kabsch coupling (Step E).} Per-sample $\mathrm{SO}(3)$-equivariant interpolation path; pre-aligns $x_0$ to $\tilde x_1$ so the field is conditioned on rotation-invariant pairs.
\item \emph{$1/(1-t)^2$ reweight (Step E delta \#1).} Activated by \texttt{target\_type=DATA} so $\mathcal{L}_\text{DATA}$ matches Proteina's $\|x_1 - \hat{x}_1\|^2 / \mathrm{nres} \cdot 1/((1-t)^2+\epsilon)$.
\item \emph{Random SO(3) aug (Step H).} Marginalises the loss over orientations; \emph{absorbed} when Kabsch is on (Step S analysis).
\item \emph{Self-conditioning (Step F).} ESMFold2 borrow \#1; warm forward at $p_\text{sc}=0.5$, zero-init $W_\text{sc}$ so a fresh model is byte-identical to non-self-cond on step 0.
\item \emph{Drop Kabsch (Step S).} Kabsch coupling at training $\neq$ from-noise sampling; removing it eliminates a known train/sample mismatch and unblocks SDE.
\item \emph{Velocity prediction (Step S).} $\hat{v}$ avoids the $(1-t)^{-1}$ blow-up that bites $\mathcal{L}_\text{DATA}$ near $t\to 1$; \texttt{target\_type=VELOCITY} drops the reweight that was double-weighting $t\to 1$ under uniform-$t$ training.
\end{itemize}

\section*{Source pointers}

\begin{itemize}\itemsep -1pt
\item Tokenizer: \texttt{src/lobster/model/latent\_generator/tokenizer/\_tokenizer\_3di\_input\_flow.py}
\item Decoder: \texttt{src/lobster/model/latent\_generator/structure\_decoder/\_vit\_decoder\_conditional.py}
\item Sampling callback: \texttt{src/lobster/model/latent\_generator/callbacks/\_flow\_backbone\_sampling.py}
\item moco interpolant: \texttt{bionemo.moco.interpolants.continuous\_time.continuous.continuous\_flow\_matching.ContinuousFlowMatcher}
\item Plan: \texttt{/cv/home/lisanzas/.cursor/plans/3di\_flow\_decoder\_01b413df.plan.md}
\end{itemize}

\end{document}
